% !TeX root = ../main.tex
\section{问题一的模型建立与求解}

\subsection{数据预处理}

\subsection{模型的建立}

% 公式用 equation 环境并编号，正文要引用它，如 式\eqref{eq:obj}。
% \begin{equation}\label{eq:obj}
%   \min_{x} \; \sum_{i=1}^{n} c_i x_i
% \end{equation}

\subsection{模型的求解}

% 算法步骤 / 伪代码。复杂度和实际运行时间要给出来。

\subsection{结果与分析}

% 给具体数字，并解释这些数字意味着什么。只贴表不解读等于没写。


\section{问题二的模型建立与求解}

\subsection{模型的建立}

\subsection{模型的求解}

\subsection{结果与分析}


\section{问题三的模型建立与求解}

\subsection{模型的建立}

\subsection{模型的求解}

\subsection{结果与分析}
